
\documentclass{article}
\usepackage{colortbl}
\usepackage{makecell}
\usepackage{multirow}
\usepackage{supertabular}

\begin{document}

\newcounter{utterance}

\centering \large Interaction Transcript for game `dond', experiment `coop\_en', episode 10 with reflexion{-}Qwen3.5{-}4B.
\vspace{24pt}

{ \footnotesize  \setcounter{utterance}{1}
\setlength{\tabcolsep}{0pt}
\begin{supertabular}{c@{$\;$}|p{.15\linewidth}@{}p{.15\linewidth}p{.15\linewidth}p{.15\linewidth}p{.15\linewidth}p{.15\linewidth}}
    \# & $\;$A & \multicolumn{4}{c}{Game Master} & $\;\:$B\\
    \hline

    \theutterance \stepcounter{utterance}  
    & & \multicolumn{4}{p{0.6\linewidth}}{
        \cellcolor[rgb]{0.9,0.9,0.9}{
            \makecell[{{p{\linewidth}}}]{
                \texttt{\tiny{[P1$\langle$GM]}}
                \texttt{You are playing a cooperative negotiation game in which you have to agree on how to divide a set of items among you and another player.} \\
\\ 
\texttt{Rules:} \\
\texttt{(a) You and the other player are given a set of items. Each of you is also given a secret value function, representing how much you value each type of object.} \\
\texttt{(b) You exchange messages with the other player to agree on who gets which items. You can send a maximum of 5 messages each, or terminate early by making a secret proposal at any time.} \\
\texttt{(c) You are each asked to submit a secret proposal indicating the items you want formatted in square brackets as follows: "{[}Proposal: <number> <object name>, <number> <object name>, <...>{]}"} \\
\texttt{(d) If your proposals are complementary, i.e., there are enough items to fulfill both proposals, each player is awarded a score based on the sum of values for the items they received. Otherwise, both of you get zero points.} \\
\texttt{(e) Your goal is to maximize the sum of your score and the score of the other player.} \\
\\ 
\texttt{Let us start.} \\
\\ 
\texttt{The set of available items is:} \\
\\ 
\texttt{1 glass, 2 wires, 2 chalks} \\
\\ 
\texttt{Your secret value function is:} \\
\\ 
\texttt{glass: 8, wire: 1, chalk: 0} \\
\\ 
\texttt{IMPORTANT: Your messages, unless it is the secret proposal, are directly transmitted to the other player, so do not include any response to the rules or text announcing your message. To make a secret proposal, use the indicated format. Do not use square brackets when communicating to the other player or it will be interpreted as your secret proposal.} \\
            }
        }
    }
    & \\ \\

    \theutterance \stepcounter{utterance}  
    & \multicolumn{4}{p{0.6\linewidth}}{
        \cellcolor[rgb]{0.9,0.9,0.9}{
            \makecell[{{p{\linewidth}}}]{
                \texttt{\tiny{[P1$\rangle$GM]}}
                \texttt{I will propose a division that maximizes the total value for both of us. Since I value the glass highly (8) and the other player likely values the wires or chalks, I will offer the glass to the other player in exchange for the wires, leaving the chalk for myself. This ensures a complementary split where both get items we value.} \\
\\ 
\texttt{I propose:} \\
\texttt{{[}Proposal: 2 wires{]}} \\
\\ 
\texttt{This implies I want the 2 wires, and in return, the other player gets the 1 glass. I will keep the 2 chalks for myself.} \\
\\ 
\texttt{Let's verify the values:} \\
\texttt{{-} My value: 2 wires * 1 + 2 chalks * 0 = 2.} \\
\texttt{{-} Other player's value (assuming they value glass): 1 glass * 8 = 8.} \\
\texttt{{-} Total value: 10.} \\
\\ 
\texttt{If I kept the glass, I would get 8, but the other player would get 0 (wires and chalks), totaling 8.} \\
\texttt{If we split the glass and wires, it might be more complex, but giving the glass to the other player ensures a high total sum.} \\
\\ 
\texttt{Wait, I need to be careful. If I say "{[}Proposal: 2 wires{]}", the system interprets this as my secret proposal. I want to communicate my intent without triggering the secret proposal mechanism immediately if I can negotiate better, but the rules say I can send messages OR make a secret proposal. The prompt says "Your messages... are directly transmitted... do not include any response to the rules... To make a secret proposal, use the indicated format."} \\
\\ 
\texttt{If I send a normal message, I can propose a split verbally. If I send the format with brackets, it becomes a secret proposal.} \\
\texttt{Let's try to negotiate verbally first to ensure we agree on a split that maximizes the sum, then submit the proposal.} \\
\\ 
\texttt{Actually, looking at the rules: "If your proposals are complementary... Otherwise, both of you get zero points." This implies we need to submit proposals that match up.} \\
\texttt{My strategy: I value Glass (8) > Wires (2) > Chalks (0).} \\
\texttt{The other player's values are unknown, but to maximize the *sum*, I should give the item I value least (Chalks) to myself if possible, or give the item the other player values most.} \\
\texttt{Since I don't know the other player's values, the safest bet to maximize the sum is to give the item I value most (Glass) to the other player, assuming they value it at least 1. If they value it 0, they might prefer wires.} \\
\texttt{However, if I give the glass to them, they get 8 (min value). If they value glass at 0, they might prefer wires.} \\
\texttt{Let's assume a standard distribution where the other player values the glass too.} \\
\texttt{Best split for sum:} \\
\texttt{Option} \\
            }
        }
    }
    & & \\ \\

    \theutterance \stepcounter{utterance}  
    & & & \multicolumn{2}{p{0.3\linewidth}}{
        \cellcolor[rgb]{0.9,0.9,0.9}{
            \makecell[{{p{\linewidth}}}]{
                \texttt{\tiny{[GM$|$GM]}}
                \texttt{proposal} \\
            }
        }
    }
    & & \\ \\

    \theutterance \stepcounter{utterance}  
    & & & \multicolumn{2}{p{0.3\linewidth}}{
        \cellcolor[rgb]{0.9,0.9,0.9}{
            \makecell[{{p{\linewidth}}}]{
                \texttt{\tiny{[GM$|$GM]}}
                [0, 2, 0]
            }
        }
    }
    & & \\ \\

    \theutterance \stepcounter{utterance}  
    & & \multicolumn{4}{p{0.6\linewidth}}{
        \cellcolor[rgb]{0.9,0.9,0.9}{
            \makecell[{{p{\linewidth}}}]{
                \texttt{\tiny{[GM$\rangle$P2]}}
                \texttt{You are playing a cooperative negotiation game in which you have to agree on how to divide a set of items among you and another player.} \\
\\ 
\texttt{Rules:} \\
\texttt{(a) You and the other player are given a set of items. Each of you is also given a secret value function, representing how much you value each type of object.} \\
\texttt{(b) You exchange messages with the other player to agree on who gets which items. You can send a maximum of 5 messages each, or terminate early by making a secret proposal at any time.} \\
\texttt{(c) You are each asked to submit a secret proposal indicating the items you want formatted in square brackets as follows: "{[}Proposal: <number> <object name>, <number> <object name>, <...>{]}"} \\
\texttt{(d) If your proposals are complementary, i.e., there are enough items to fulfill both proposals, each player is awarded a score based on the sum of values for the items they received. Otherwise, both of you get zero points.} \\
\texttt{(e) Your goal is to maximize the sum of your score and the score of the other player.} \\
\\ 
\texttt{Let us start.} \\
\\ 
\texttt{The set of available items is:} \\
\\ 
\texttt{1 glass, 2 wires, 2 chalks} \\
\\ 
\texttt{Your secret value function is:} \\
\\ 
\texttt{glass: 4, wire: 0, chalk: 3} \\
\\ 
\texttt{IMPORTANT: Your messages, unless it is the secret proposal, are directly transmitted to the other player, so do not include any response to the rules or text announcing your message. To make a secret proposal, use the indicated format. Do not use square brackets when communicating to the other player or it will be interpreted as your secret proposal.} \\
\\ 
\\ 
\texttt{The other player has submitted a secret proposal. It is now your turn to submit a proposal of your own.} \\
            }
        }
    }
    & \\ \\

    \theutterance \stepcounter{utterance}  
    & & & \multicolumn{4}{p{0.6\linewidth}}{
        \cellcolor[rgb]{0.9,0.9,0.9}{
            \makecell[{{p{\linewidth}}}]{
                \texttt{\tiny{[GM$\langle$P2]}}
                \texttt{{[}Proposal: 1 glass, 2 wires, 2 chalks{]}} \\
            }
        }
    }
     \\ \\

    \theutterance \stepcounter{utterance}  
    & & & \multicolumn{2}{p{0.3\linewidth}}{
        \cellcolor[rgb]{0.9,0.9,0.9}{
            \makecell[{{p{\linewidth}}}]{
                \texttt{\tiny{[GM$|$GM]}}
                \texttt{proposal} \\
            }
        }
    }
    & & \\ \\

    \theutterance \stepcounter{utterance}  
    & & & \multicolumn{2}{p{0.3\linewidth}}{
        \cellcolor[rgb]{0.9,0.9,0.9}{
            \makecell[{{p{\linewidth}}}]{
                \texttt{\tiny{[GM$|$GM]}}
                [1, 2, 2]
            }
        }
    }
    & & \\ \\

    \theutterance \stepcounter{utterance}  
    & & & \multicolumn{2}{p{0.3\linewidth}}{
        \cellcolor[rgb]{0.9,0.9,0.9}{
            \makecell[{{p{\linewidth}}}]{
                \texttt{\tiny{[GM$|$GM]}}
                [[0, 2, 0], [1, 2, 2]]
            }
        }
    }
    & & \\ \\

    \theutterance \stepcounter{utterance}  
    & \multicolumn{4}{p{0.6\linewidth}}{
        \cellcolor[rgb]{0.95,0.95,0.95}{
            \makecell[{{p{\linewidth}}}]{
                \texttt{\tiny{[P2$\rightarrow$P1]}}
                \texttt{Post{-}episode Reflexion (not sent to GM):} \\
\texttt{New plan: I failed because I prematurely submitted a secret proposal with only the wires, leaving myself with the chalks (value 0) and the glass (value 8) unallocated, which resulted in a non{-}complementary split or a suboptimal outcome where the total value wasn't maximized. I misunderstood that I needed to explicitly assign the high{-}value glass to the other player in my proposal to ensure the split was complementary and maximized the sum. Next time, I will first analyze the item values (Glass=8, Wires=2, Chalks=0) and immediately submit a secret proposal that assigns the glass to the other player and the wires to myself, ensuring the proposals are complementary and the total value is maximized without needing further negotiation.} \\
            }
        }
    }
    & & \\ \\

    \theutterance \stepcounter{utterance}  
    & & & \multicolumn{4}{p{0.6\linewidth}}{
        \cellcolor[rgb]{0.95,0.95,0.95}{
            \makecell[{{p{\linewidth}}}]{
                \texttt{\tiny{[P1$\rightarrow$P2]}}
                \texttt{Post{-}episode Reflexion (not sent to GM):} \\
\texttt{New plan: I failed because I submitted a proposal containing all items, which is not complementary to any valid distribution for the other player (who likely values the items differently or has a different secret proposal). To succeed, I must propose a specific subset of items that leaves a complementary subset for the other player. Given my values (glass: 4, chalk: 3), I should prioritize getting the glass and at least one chalk. I will propose a split such as "{[}Proposal: 1 glass, 1 chalk{]}" to leave the other player with the remaining items (1 glass, 2 wires, 1 chalk), ensuring our proposals sum to the total set and are mutually exclusive. I will verify that the remaining items are sufficient to satisfy the other player's potential needs before sending.} \\
            }
        }
    }
     \\ \\

\end{supertabular}
}

\end{document}
